\chapter{Radiofrequency Catheter Ablation}

\section*{Introduction \& Clinical Indications}
Radiofrequency catheter ablation uses controlled thermal energy to create small areas of scar that interrupt an abnormal electrical circuit or isolate an arrhythmogenic focus. It is used for symptomatic or clinically important arrhythmias such as supraventricular tachycardia, atrial flutter, selected atrial fibrillation, and some ventricular arrhythmias. Cryothermal energy and other technologies may be used in selected cases, but the procedural principles and specialist setting differ.

The decision to ablate considers symptom burden, arrhythmia type, structural heart disease, stroke risk, medication response, patient preference, and procedural risk. Catheter ablation does not automatically eliminate the need for anticoagulation, especially after atrial-fibrillation treatment; this depends on thromboembolic risk and specialist guidance.

\section*{Relevant Surgical Anatomy}
Targets are defined by the cardiac conduction system and the mechanism of the arrhythmia. The atrioventricular node, His bundle, coronary arteries, phrenic nerves, pulmonary veins, left atrial appendage, and esophagus may be near ablation sites. Venous access leads to the right heart; transseptal puncture may be required for left-atrial procedures.

\section*{Preoperative Workup \& Preparation}
Obtain ECG documentation, rhythm monitoring, echocardiography when indicated, medication review, renal function, anticoagulation assessment, and imaging for selected left-atrial procedures. Explain vascular injury, bleeding, stroke, transient ischemic attack, cardiac perforation, tamponade, heart block, phrenic-nerve injury, pulmonary-vein stenosis, esophageal injury, recurrence, need for a pacemaker, and repeat treatment. Anticoagulation and antiarrhythmic management follow the specific procedure and patient risk.

\section*{Surgical Technique \& Procedure}
Under local anesthesia with sedation or general anesthesia, catheters are introduced through venous access and positioned in the heart. Intracardiac recordings and three-dimensional mapping identify the arrhythmogenic substrate. Radiofrequency energy is delivered at selected points while temperature, power, impedance, and neighboring structures are monitored. The endpoint may be termination of the arrhythmia, bidirectional block across a circuit, or electrical isolation of a target region. Observation and repeat testing assess immediate success.

\section*{Complications \& Risks}
Risks include groin bleeding or hematoma, vascular injury, stroke, transient ischemic attack, cardiac perforation, tamponade, coronary injury, atrioventricular block, phrenic-nerve injury, pulmonary-vein stenosis, esophageal injury, infection, recurrence, and death. Chest pain, severe breathlessness, fainting, neurologic symptoms, fever, difficulty swallowing, or uncontrolled bleeding requires urgent assessment.

\section*{Postoperative Care \& Recovery}
Monitor the access sites, rhythm, blood pressure, neurologic status, and symptoms. Provide medication and anticoagulation instructions, activity restrictions, rhythm follow-up, and wound-care advice. Early transient palpitations do not always indicate failure, but sustained or severe symptoms require assessment. Long-term rhythm monitoring and repeat ablation may be needed.

\section*{Outcomes \& Prognosis}
Ablation can markedly reduce arrhythmia burden and improve symptoms, but success and recurrence rates depend on the arrhythmia, heart structure, comorbidities, and prior treatment. Ongoing management of hypertension, sleep apnea, obesity, alcohol exposure, and other risk factors improves long-term rhythm outcomes.

\section*{References}
\begin{enumerate}
\item Heart Rhythm Society/European Heart Rhythm Association. Expert consensus on catheter and surgical ablation of atrial fibrillation.
\item European Society of Cardiology. Guidelines for supraventricular tachycardia and atrial fibrillation.
\item American College of Cardiology/American Heart Association/Heart Rhythm Society. Guideline for the Management of Patients With Ventricular Arrhythmias.
\end{enumerate}
